%================================
% Vocabulary & Language Focus — XB3 (Unit 1)
% 编译方式: xelatex main.tex (需运行两遍以生成目录)
%================================
\documentclass[../main.tex]{subfiles}

\begin{document}

\section{Unit 1}

\vw{blooming}{/ˈbluːmɪŋ/}{adj.}

\sense{adj.}{盛开的，开花的}
\col{blooming flowers/trees/gardens}
\ex{The blooming cherry trees lined the avenue like a corridor of pink clouds.}
\ex{She picked a small bunch of blooming daisies from the meadow.}

\sense{adj.}{青春焕发的；容光焕发的（拓展）}
\col{a blooming bride/youth}
\ex{She looked a blooming nineteen-year-old on her wedding day.}

\begin{origin}
直接借自古英语 \textit{blōmian}（开花、绽放），名词 bloom 来自原始日耳曼语 *\textit{blōmaz}，一说与 \textit{blōwan}（绽放、繁荣）同根，可能与"闪光、明亮"的概念相关（花之色即花之耀）。同源词 \textbf{bloom}（花；青春；焕发）、\textbf{blossom}（果树花；盛开）——blossom 偏指果树（apple blossom 苹果花），bloom 可用于一切花。
\end{origin}

\vw{unfold}{/ʌnˈfəʊld/}{vi. \& vt.}

\sense{vt. \& vi.}{展开，打开}
\col{unfold a map/newspaper/letter}
\ex{She unfolded the map and spread it across the table.}
\col{unfold one's arms/wings}
\ex{The eagle unfolded its wings and rose from the cliff.}

\sense{vi.}{开花}
\ex{The roses began to unfold in the warm spring sunshine.}

\sense{vi.}{展现，披露}
\col{a story/scene unfolds}
\ex{As the plot unfolds, we learn that the hero is not who he seems.}
\col{unfold before one's eyes}
\ex{The whole valley unfolded before us from the summit.}

\begin{origin}
古英语 \textit{unfealdan}：\textit{un-}（否定前缀$\to$反向）+ \textit{fealdan}（折叠）。字面义即"反向折叠$\to$展开"。同源词 \textbf{fold}（折叠）、\textbf{folder}（文件夹）、\textbf{folderol}（废话）、\textbf{faldage}（圈羊栏）。注意 unfold 的"展现、披露"义是隐喻——"折拢的故事一页页打开"，是叙事学里"情节铺陈"的标准动词。
\end{origin}

\vw{ethnically}{/ˈeθnɪkli/}{adv.}

\sense{adv.}{种族地；人种上}
\col{ethnically diverse/mixed/cleansed}
\ex{The neighbourhood is one of the most ethnically diverse in the city.}
\col{ethnically distinct peoples}
\ex{The country is home to several ethnically distinct peoples, each with its own language and customs.}

\begin{lex}{ethnic / racial / tribal}
\textbf{ethnic} 偏文化语言学意味（族群：语言、宗教、传统），是中性偏正式用词；\textbf{racial} 严格指生物学上的"人种"（race），因 20 世纪历史创伤而承载沉重语义，当代敏感写作倾向以 ethnic 替代；\textbf{tribal} 指部落社会结构（非洲、南美原住民），不可与前两者互换。
\end{lex}

\begin{origin}
借自希腊语 \textit{ethnos}（民族、人民），原始印欧语 *\textit{swe-dh-}（自己的、本族的——与 self, swain 同根）。同源词 \textbf{ethnology}（民族学）、\textbf{ethnography}（民族志）、\textbf{ethnocentric}（民族中心主义的）。希腊城邦时代 ethnos 用以指"非希腊人"，带有他者化色彩，19 世纪才被人类学重新中立化。
\end{origin}

\vw{distinct}{/dɪˈstɪŋkt/}{adj.}

\sense{adj.}{截然不同的，有区别的}
\col{distinct from sth}
\ex{The two species are quite distinct from each other in both colour and song.}
\col{a distinct possibility/advantage}
\ex{There is a distinct possibility that the meeting will be postponed.}

\sense{adj.}{不同种类的}
\col{two/three distinct groups/types}
\ex{The book identifies three distinct types of climate in the region.}

\sense{adj.}{清楚的，清晰的}
\col{a distinct memory/impression/voice}
\ex{He has a distinct memory of the day his grandfather took him fishing.}
\col{a distinct smell/taste/sound}
\ex{There was a distinct smell of smoke coming from the kitchen.}

\begin{gram}{distinct / distinctive / distinguishable}
\textbf{distinct} 既可表"不同的"也可表"清晰的"（\textit{two distinct ideas}; \textit{a distinct image}）；\textbf{distinctive} 只能表"有特色的"（\textit{a distinctive accent}），不可表"清晰的"；\textbf{distinguishable} 表"可辨别的"（多用于否定：\textit{The twins are barely distinguishable}）。三词易混。
\end{gram}

\begin{origin}
借自拉丁语 \textit{distinctus}（分开、区别开），是 \textit{distingere}（区分）的过去分词：\textit{dis-}（分开）+ \textit{stinguere}（刺、点——原始印欧语 *\textit{steig-} 刺）。同源词 \textbf{distinguish}（区分）、\textbf{distinguished}（杰出的）、\textbf{extinct}（熄灭，"火苗被分开"）、\textbf{instinct}（本能，"从内部被刺入的冲动"）。
\end{origin}

\vw{territory}{/ˈterətri/}{n.}

\sense{n.}{（加拿大）地区}
\ex{The Yukon is a federal territory in the northwest of Canada.}

\sense{n.}{领土，领地}
\col{territorial waters/airspace}
\ex{Chinese fishing vessels were found inside Indonesian territorial waters.}
\col{occupied/disputed territory}
\ex{The two countries have fought over the disputed territory for decades.}

\sense{n.}{领域，范围}
\col{territory for sb/sth}
\ex{This question falls outside my territory — you'd better ask a lawyer.}
\col{uncharted territory 未知领域}
\ex{The team is moving into uncharted territory with its new project.}

\begin{origin}
借自拉丁语 \textit{territorium}（一片土地、辖区），源于 \textit{terra}（土地）。同源词 \textbf{terrace}（梯田；阳台；露天看台）、\textbf{terrain}（地形）、\textbf{terrestrial}（陆地的；地球的）、\textbf{Mediterranean}（medius 中间 + terra 土地"$\to$陆中之海"）、\textbf{territory} 与 \textbf{exterior} 是 *\textit{ters-}（干、燥）原始印欧词根的远亲——"陆地"本是"干燥之处"，与"海"相对。
\end{origin}

\vw{valley}{/ˈvæli/}{n.}

\sense{n.}{山谷，溪谷}
\col{a river/green/glacial valley}
\ex{The village sits in a deep green valley surrounded by pine forests.}
\col{the valley floor 谷底}
\ex{Flood waters rushed down the valley floor after the storm.}

\begin{origin}
古英语 \textit{feallan}（落下）派生的 \textit{valley}，本义"位于低洼处的地"——与 fall 同源，valley 即"低处"。同源词 \textbf{vale}（诗歌用语"谷"：\textit{ vale of tears} 泪之谷，生之苦海）、\textbf{valley}（日常用语）、\textbf{dismal}（dismal 阴郁的，本是拉丁语"双十"$\to$倒霉日）。
\end{origin}

\vw{carve}{/kɑːv/}{vt.}

\sense{vt.}{侵蚀，冲刷}
\col{carve sth out（河流、风等）}
\ex{Over centuries the river has carved out a deep gorge through the limestone.}
\col{carve through/across sth}
\ex{The path carved its way across the mountainside.}

\sense{vt.}{雕，刻}
\col{carve sth (in/into/on sth)}
\ex{The artist carved a delicate pattern into the wood.}
\col{carve a statue/tomb/inscription}
\ex{The tomb was carved with figures of warriors and horses.}

\sense{vt.}{切（肉）}
\col{carve the turkey/roast}
\ex{Will you carve the turkey, or shall I?}

\begin{gram}{carve 的"塑造"隐喻丛}
\textbf{carve out a career/niche/reputation}（开辟事业/职位/名声），\textbf{carve up}（瓜分；切碎——多含贬义：\textit{The victors carved up the defeated country}），\textbf{carve one's way through}（开辟道路）。源于雕刻匠的"凿出形状"，与 sculpt 形异义近。注意 carve out a niche 是当代职场高频搭配："在市场中凿出一块专属空间"。
\end{gram}

\begin{origin}
古英语 \textit{ceorfan}（切割、雕刻），原始日耳曼语 *\textit{kerfan}。同源词 \textbf{kerf}（锯缝）、古英语 \textit{ceorfan} 派生的 \textbf{scarf}（围巾；嵌接；切口）。希腊语 \textit{grapho}（写）和 carve 是不同词族，"刻"与"写"在英语中分家——但德语 \textit{kerben} 仍保留"切刻"义。
\end{origin}

\vw{mighty}{/ˈmaɪti/}{adj.}

\sense{adj.}{巨大的，非凡的}
\col{a mighty river/mountain/oak}
\ex{We stood on the bank of the mighty Mississippi at sunset.}
\col{mighty effort/war}
\ex{The battle was a mighty struggle that changed the course of history.}

\sense{adj.}{强而有力的}
\col{a mighty ruler/nation/warrior}
\ex{The Roman Empire was once the mightiest nation in the western world.}

\begin{lex}{mighty / powerful / strong / forceful}
\textbf{mighty} 偏文学/庄严，描述压倒性力量（\textit{mighty deeds} 丰功伟绩）；\textbf{powerful} 通用，能力/影响力大（\textit{a powerful engine} 强力引擎；\textit{powerful speech} 有力演讲）；\textbf{strong} 强调耐久/抵抗力（\textit{strong tea} 浓茶；\textit{a strong personality} 坚强个性）；\textbf{forceful} 强调推动/推动力（\textit{a forceful argument} 有说服力的论据）。美口语中 \textbf{mighty} 还可作副词表"非常"（\textit{mighty tired}），属旧用法。
\end{lex}

\begin{origin}
古英语 \textit{mihtig}（强大的），是 \textit{mæht}（力量、能力）的形容词形式，原始日耳曼语 *\textit{mahtiz}。同源词 \textbf{might}（力量；\textit{might and main} 竭尽全力）、\textbf{main}（主要的；体力）。may（可能）和 might 同根——"可能"在古人看来即是"潜在力量"。
\end{origin}

\vw{splendid}{/ˈsplendɪd/}{adj.}

\sense{adj.}{壮丽的，雄伟的}
\col{a splendid palace/view/sunset}
\ex{From the hilltop we enjoyed a splendid view of the snow-capped mountains.}
\ex{The Taj Mahal is one of the most splendid monuments in the world.}

\sense{adj.}{极佳的，非常好的}
\col{a splendid idea/performance/holiday}
\ex{We had a splendid time at the beach last summer.}
\col{splendid isolation 光荣孤立（19 世纪英国外交政策）}
\ex{After the fall of Napoleon, Britain returned to its policy of splendid isolation.}

\begin{origin}
17 世纪借自拉丁语 \textit{splendidus}（光辉的、华丽的），源于 \textit{splendere}（发光、闪耀）。同源词 \textbf{splendour}（壮丽；辉煌）、\textbf{resplendent}（光辉灿烂的：re- 加强 + splendere）、\textbf{resplendence}（灿烂）。原始印欧语 *\textit{splend-}（发光）一说与希腊语 \textit{aspalathos}（香料植物）有关——但此说存疑。splendid 与 splendidly 都属褒义强度高的书面用词。
\end{origin}

\vp{apart from}

\sense{短语}{除……外（还），此外}
\ex{Apart from English, she also speaks fluent French and German.}
\ex{Apart from the high cost, the plan has other disadvantages.}

\sense{短语}{除……外（都），要不是}
\ex{Apart from a few scratches, the car was undamaged.}
\ex{Apart from a brief shower, the weather was perfect all week.}

\begin{lex}{apart from / besides / except / in addition to}
\textbf{apart from} 在英式英语中常表"除……外还"或"除……外都"，含义需依上下文判定；\textbf{besides} 同 apart from 的第一义（"还"），不可表"除……外都"；\textbf{except} 仅表"除……外都"，后不接"还"；\textbf{in addition to} 仅表"除……外还"。美式英语中 apart from 多半等同 except。一句话：\textit{Apart from / Besides English, he speaks French.}（"还会说法语"——肯定）。
\end{lex}

\vw{urban}{/ˈɜːbən/}{adj.}

\sense{adj.}{城市的，城镇的}
\col{urban areas/life/centres}
\ex{More than half the world's population now lives in urban areas.}
\col{urban planning/renewal}
\ex{The mayor unveiled a new urban-renewal plan for the downtown area.}

\sense{adj.}{具有都市风格的（拓展）}
\col{urban sophisticates}
\ex{The magazine is aimed at young, urban professionals.}

\begin{origin}
17 世纪借自拉丁语 \textit{urbanus}（城市的；彬彬有礼的），源于 \textit{urbs}（城市）。同源词 \textbf{urbane}（有礼貌的、文雅的——源自"城里人"比"乡巴佬"更礼貌的偏见）、\textbf{suburb}（sub- 下 + urbs $\to$城下的居所）、\textbf{conurbation}（集合城市）。"城市"和"礼貌"在英语里共享词根，汉语"礼"与"里"无关，乃不同文化之趣。
\end{origin}

\vw{ancestor}{/ˈænsestə(r)/}{n.}

\sense{n.}{祖宗，祖先}
\col{ancestors of sb}
\ex{My ancestors came from a small village in the north of England.}
\col{ancestor worship 祖先崇拜}
\ex{Ancestor worship is a central feature of many traditional cultures.}

\sense{n.}{原种；原型}
\col{ancestor of sth}
\ex{The giant panda is thought to be a descendant of an ancient bear-like ancestor.}
\ex{The modern printing press is a direct ancestor of today's digital devices.}

\begin{origin}
13 世纪借自古法语 \textit{ancestre}（祖宗），晚期拉丁语 \textit{antecessor}（先行者）：\textit{ante-}（在前）+ \textit{cedere}（走、行）。同源词 \textbf{antecede}（先行）、\textbf{antecedent}（前事；先行词；祖先）、\textbf{ancestry}（世系；血统）。注意词形是 -cestor 而非 -cester，是古法语内部音变所致。
\end{origin}

\vw{indigenous}{/ɪnˈdɪdʒənəs/}{adj.}

\sense{adj.}{本地的，土生土长的}
\col{indigenous to sth}
\ex{Koalas are indigenous to Australia and are not found naturally anywhere else.}
\col{indigenous peoples/communities}
\ex{The conference discussed the rights of indigenous peoples in the Arctic region.}

\begin{lex}{indigenous / native / aboriginal / endemic}
\textbf{indigenous} 正式语体，人类学/地理学/生态学首选；\textbf{native} 日常英语通用，亦可指"出生地的"（\textit{her native country} 她的祖国）；\textbf{aboriginal} 多指澳大利亚原住民（\textit{Aboriginal Australians}）或史前人类（\textit{aboriginal inhabitants}），日常少用；\textbf{endemic} 严格限于生态学/医学（\textit{an endemic species} 特有种；\textit{an endemic disease} 地方病）。
\end{lex}

\begin{origin}
17 世纪借自拉丁语 \textit{indigena}（本地的、当地的），由 \textit{indi-}（在……内；变体于 \textit{indu-}，与 \textit{in} 同源）+ \textit{gignere}（生育）构成——本义"国内出生的"。同源词 \textbf{indignant}（愤怒的；字面"不配的"）、\textbf{indignity}（侮辱）、\textbf{genuine}（真正的；"天生的"）、\textbf{generate}（产生）。
\end{origin}

\vw{strait}{/streɪt/}{n.}

\sense{n.}{海峡，水道}
\col{the Strait(s) of ...}
\ex{The ship sailed through the Strait of Gibraltar into the Mediterranean.}
\col{international straits 国际海峡}
\ex{Freedom of navigation through international straits is guaranteed by international law.}

\begin{origin}
古法语 \textit{estreit}（狭窄的），拉丁语 \textit{strictus}（拉紧的、狭窄的），是 \textit{stringere}（拉紧）的过去分词。同源词 \textbf{strict}（严格的）、\textbf{straitjacket}（约束衣）、\textbf{constrain}（强制）、\textbf{restrain}（抑制）、\textbf{stringent}（严格的、紧缩的）。海峡之"狭"与限制之"严"是同一意象。
\end{origin}

\vw{minority}{/maɪˈnɒrəti/}{n.}

\sense{n.}{少数民族，少数群体}
\col{a minority group/ethnic/religious}
\ex{The law protects the rights of ethnic and religious minorities.}
\col{minority rights/protection}
\ex{The new bill aims to strengthen minority rights in the workplace.}

\sense{n.}{少数，少数人}
\col{a/the minority of}
\ex{Only a minority of voters supported the proposal.}
\col{be in the minority 占少数}
\ex{At the meeting, women were very much in the minority.}

\begin{origin}
16 世纪借自法语 \textit{minorité} $\leftarrow$ 中世纪拉丁语 \textit{minoritas}，源于拉丁语 \textit{minor}（更小的，比较级）。同源词 \textbf{minor}（较小的；次要的；未成年人）、\textbf{minus}（减去；零下）、\textbf{minuscule}（极小的）、\textbf{minister}（伺候较小者$\to$侍从、部长）、\textbf{administer}（管理）。minority 形上即"少数派"，majority（多数派）是其镜像。
\end{origin}

\vw{cuisine}{/kwɪˈziːn/}{n.}

\sense{n.}{烹饪，风味}
\col{French/Italian/Chinese cuisine}
\ex{Paris is famous for its fine cuisine and excellent restaurants.}
\col{local/regional cuisine}
\ex{Every region of Italy has its own distinctive cuisine.}

\sense{n.}{饭菜，菜肴}
\col{a cuisine of ...}
\ex{The hotel offers an excellent cuisine of fresh local dishes.}

\begin{origin}
19 世纪借自法语 \textit{cuisine}（厨房；烹饪），拉丁语 \textit{coquina}（烹饪的艺术），源于 \textit{coquere}（烹、煮）。同源词 \textbf{kitchen}（厨房）、\textbf{cook}（厨师；烹调）、\textbf{concoct}（调制、编造）、\textbf{precarious}（不稳定的，本义"经祈祷求得的"——与 cook 无关）。cuisine 严格说是"厨房或烹饪"，现代英语偏指后者。
\end{origin}

\vw{dynamic}{/daɪˈnæmɪk/}{adj.}

\sense{adj.}{（事物）有活力的}
\col{a dynamic person/city/economy}
\ex{Hong Kong is one of the most dynamic cities in Asia.}
\col{dynamic leadership}
\ex{The company is looking for a dynamic leader to drive its next stage of growth.}

\sense{adj.}{动态的，发展变化的}
\col{dynamic process/equilibrium/systems}
\ex{The market is in a state of dynamic equilibrium, with prices constantly adjusting.}

\sense{adj.}{充满活力的}
\col{a dynamic personality/performance}
\ex{Her dynamic personality made her a natural choice for the role of team leader.}

\begin{gram}{dynamic / dynamics 的词性差异}
\textbf{dynamic}（形容词）"有活力的"；\textbf{dynamics}（名词，单数或复数均可）"动力学"（物理：\textit{fluid dynamics} 流体力学）或"动力、势态"（社会：\textit{the dynamics of group behaviour} 群体行为的动力机制）。注意作"动力"解时是不可数集合名词。
\end{gram}

\begin{origin}
19 世纪借自法语 \textit{dynamique}，希腊语 \textit{dynamikos}（有力的），源于 \textit{dynamis}（力量、能力）。同源词 \textbf{dynamite}（炸药——诺贝尔 1867 年用希腊词 dynamis 命名）、\textbf{dynasty}（王朝，统治者的力量$\to$世系）、\textbf{ dynamo}（发电机）、\textbf{odometer}（里程表——hodos 路径，与 dyn- 不同源）。dynam- 词族借自希腊后多与"力"相关。
\end{origin}

\vw{sample}{/ˈsɑːmpl/}{vt. \& n.}

\sense{vt.}{尝，品尝，体验}
\col{sample sth}
\ex{I'd like to sample some of the local wines before buying.}
\ex{Visitors to the theme park can sample the rides for free on opening day.}

\sense{vt.}{抽样检验}
\col{sample sth (for sth)}
\ex{Soil from the site was sampled for traces of contaminants.}

\sense{n.}{样本，样品}
\col{a sample of sth}
\ex{The shop gives away free samples of its new perfume.}
\col{a random/randomised sample 随机样本}
\ex{Polling is based on a random sample of voters.}

\sense{n.}{取样，抽样}
\col{sample size/sampling 样本量/抽样}
\ex{A larger sample size will give more reliable results.}

\begin{origin}
13 世纪起作"范例、样品"义，借自古法语 \textit{essample}（例子）$\leftarrow$ 拉丁语 \textit{exemplum}（例子）。动词"抽样、尝"由名词义"样品"回溯引申。同源词 \textbf{example}（例子）、\textbf{exemplar}（典范）、\textbf{exemplify}（例证）。sample 与 example 形义皆近，区别在于 sample 偏具体/物质（尝、看、用），example 偏抽象/说明（举出、示范）。
\end{origin}

\vp{off the beaten track}

\sense{短语}{不落俗套，打破常规}
\ex{Her latest novel is delightfully off the beaten track — a satire on the publishing world.}

\sense{短语}{远离闹市，偏远}
\ex{We spent our holiday off the beaten track in a remote mountain village.}
\ex{For travellers looking to get off the beaten track, the west coast of the island is ideal.}

\begin{origin}
源于登山/徒步文化：beaten track 即"被脚步踩出的路"——人迹罕至处便是"被打扰最少"的地方。同构习语 \textbf{beaten path}（熟路；惯例）、\textbf{go off the rails}（出轨；行为失常）、\textbf{on the right track}（方向正确）。track/race/path/trail 皆印欧 *\textit{trek-}（拉、拖）之裔。
\end{origin}

\vp{come as no surprise}

\sense{短语}{不足为奇，毫不奇怪}
\col{It comes as no surprise that ...}
\ex{It comes as no surprise that the company has been losing money — the management ignored the warning signs.}
\ex{That he won the prize came as no surprise to anyone who knew his work.}

\begin{gram}{形式主语 it + 被动语态 + that 从句}
\textbf{It comes as no surprise that...} 是英文"延迟主语"（extraposition）的典型结构：that 从句本可作主语（\textit{That he won the prize came as no surprise.}），但因从句过长或为避免头重脚轻，移至句尾而由 it 形式占据主语位。相似结构：\textit{It is well known that...}、\textit{It turns out that...}、\textit{It occurred to me that...}。写作中频率极高。
\end{gram}

\vw{depth}{/depθ/}{n.}

\sense{n.}{最深处}
\col{at a depth of ...}
\ex{The submarine was operating at a depth of two hundred metres.}

\sense{n.}{深度}
\col{in depth 深入地}
\ex{We need to discuss the issue in greater depth before deciding.}
\col{the depth of sth}
\ex{The depth of the lake varies from a few metres to over thirty.}

\sense{n.}{深刻；深厚}
\col{depth of feeling/knowledge/understanding}
\ex{She writes with great depth of feeling about her childhood in the mountains.}
\col{be out of one's depth 力不从心}
\ex{I felt completely out of my depth during the advanced seminar.}

\begin{gram}{the depths of ...：强烈情感"深处"的搭配}
\textbf{in the depths of despair/winter/the countryside/the ocean}。depth 的复数常承载感情色彩："在最深的绝望中"（\textit{in the depths of despair}）、"在最深的冬日里"（\textit{in the depths of winter}）。比单纯 in despair 文学性更强。
\end{gram}

\begin{origin}
古英语 \textit{deop}（深的）派生的名词 \textit{deopþ}，原始日耳曼语 *\textit{deupaz}。同源词 \textbf{deep}（深的）、\textbf{deeply}（深深地）、\textbf{deepen}（加深）、古英语 \textit{dīpan}（浸）。deep-depth 一族与"水"无直接词源关系——它们的祖义是"向下、远处"，与 dip 远亲。
\end{origin}

\vw{botanical}{/bəˈtænɪkl/}{adj.}

\sense{adj.}{植物学的}
\col{botanical garden/name}
\ex{We spent a sunny afternoon in the botanical garden, admiring the rare orchids.}
\ex{Every plant in the garden has a small botanical name card beside it.}

\begin{origin}
17 世纪借自拉丁语 \textit{botanicus}，希腊语 \textit{botanikos}（关于草本的），源于 \textit{botanē}（草、植物）、\textit{boskein}（放牧吃草）。同源词 \textbf{botany}（植物学）、\textbf{botanist}（植物学家）、\textbf{botanize}（采集研究植物）。希腊语中"草"与"放牧"同源——植物学本就是"研究羊吃什么"的学问。
\end{origin}

\vw{blanket}{/ˈblæŋkɪt/}{n.}

\sense{n.}{厚层}
\col{a blanket of sth}
\ex{A thick blanket of snow covered the fields overnight.}
\ex{The valley was hidden under a blanket of mist at dawn.}

\sense{n.}{毛毯}
\col{a woollen/cotton blanket}
\ex{Her grandmother wrapped her in a soft woollen blanket by the fire.}

\begin{gram}{blanket 作形容词的"全覆盖"用法}
\textbf{a blanket ban/term/approach/coverage}（全面禁止／通用术语／一刀切的方法／全面报道）。blanket 由名词（覆盖一切之"毯"）转化为形容词，修饰抽象政策/方法时强调"无差别覆盖"。例：\textit{The government imposed a blanket ban on single-use plastics.}
\end{gram}

\begin{origin}
13 世纪借自古法语 \textit{blanchet}（小块白布），由 \textit{blanc}（白）派生——blanket 本是"白色羊毛小毯"。同源词 \textbf{blancmange}（白色奶冻）、\textbf{blank}（空白的；\textit{blanket} 形近不同源）。注意 blanket 原色为白，今天各式花色皆有。
\end{origin}

\vw{welfare}{/ˈwelfeə(r)/}{n.}

\sense{n.}{幸福，福祉}
\col{the welfare of sb/sth}
\ex{The company is committed to the welfare of its employees.}
\col{child/animal welfare 儿童福利/动物福利}
\ex{She has spent her career campaigning for animal welfare.}

\sense{n.}{福利}
\col{social/welfare system/state}
\ex{Scandinavian countries are well known for their generous welfare systems.}
\col{on welfare 接受福利救济}
\ex{After losing his job, the family had to go on welfare for several months.}

\begin{origin}
古英语 \textit{wel}（好）+\textit{faren}（行进、生活），本义"生活得好"——"福利"是远引申。同源词 \textbf{well}（好；井）、\textbf{fare}（进展；车费）、\textbf{fair}（公平的；集市，与 fare 形近但实为拉丁语 \textit{feria} 节日）、\textbf{unwell}（不适）。welfare 与 well-being（安康、福祉）部分重叠，但 welfare 多含"政府救济"义。
\end{origin}

\vw{accommodation}{/əˌkɒməˈdeɪʃn/}{n.}

\sense{n.}{住处，膳宿}
\col{hotel/living/luxury accommodation}
\ex{The price of the tour includes flights and hotel accommodation.}
\col{rented/self-catered accommodation}
\ex{Accommodation is hard to find in the city during the festival week.}

\begin{gram}{英美拼写差异：accommodation / accommodations}
英式英语 \textbf{accommodation} 作不可数名词（\textit{accommodation is included}）；美式英语常用复数 \textbf{accommodations}（\textit{the accommodations are excellent}）。同类的英美拼写差异：colour/color, theatre/theater, programme/program。
\end{gram}

\begin{origin}
16 世纪借自拉丁语 \textit{accommodatio}（适应、迁就），源于 \textit{accommodare}（使适应）：\textit{ac-}（ad- 的同化形式"朝向"）+ \textit{commodus}（合适的）。同源词 \textbf{accommodate}（容纳；适应；迁就）、\textbf{commodious}（宽敞的）、\textbf{accommodating}（随和的、好商量的）。住处之"accommodation"是"使旅人得到迁就/安置"的语义延伸。
\end{origin}

\vw{harmony}{/ˈhɑːməni/}{n.}

\sense{n.}{融洽，和睦}
\col{live/work in harmony}
\ex{The two communities have learned to live in harmony with one another.}
\col{domestic harmony 家庭和睦}
\ex{The traditional view places domestic harmony above individual freedom.}

\sense{n.}{和谐}
\col{colour/musical harmony 色彩和谐/音乐和声}
\ex{The painting is a study in greens and blues in perfect harmony.}
\col{in harmony with sth 与……和谐一致}
\ex{Our new building is designed to be in harmony with the surrounding landscape.}

\begin{lex}{harmony / agreement / unity}
\textbf{harmony} 强调元素之间的相容性、和声感，多指音乐/色彩/关系；\textbf{agreement} 强调意见一致，多指合约/共识；\textbf{unity} 强调凝聚为一个整体，常含政治/国家意涵。family harmony（家庭和睦）不可说成 family agreement——后者指"达成协议"。
\end{lex}

\begin{origin}
14 世纪借自古法语 \textit{armonie}，拉丁语 \textit{harmonia}，希腊语 \textit{harmonia}（契合、连接），源于 \textit{harmos}（关节、接合）。同源词 \textbf{harmonic}（和声的；谐波的）、\textbf{harmonica}（口琴）、\textbf{harmonize}（协调；配和声）、\textbf{harmonium}（簧风琴）。希腊神话中 Harmonia 是战神与爱神结合所生——和声源自战火与爱情的结合。
\end{origin}

\vw{sustain}{/səˈsteɪn/}{vt.}

\sense{vt.}{使保持，使稳定持续}
\col{sustain interest/growth/momentum}
\ex{The country has sustained rapid economic growth for over a decade.}
\col{sustain a lead/relationship}
\ex{The team managed to sustain their lead until the final whistle.}

\sense{vt.}{维持（生命、生活）}
\col{sustain sb on sth}
\ex{The survivors were sustained on a diet of fish and wild berries.}
\col{sustain life/living}
\ex{The soil is too poor to sustain life without irrigation.}

\sense{vt.}{承受，遭受（损伤）}
\col{sustain damage/injuries/losses}
\ex{The building sustained serious damage in the earthquake.}

\begin{gram}{sustain / maintain / support / uphold}
\textbf{sustain} 强调"持续不断"地保持（生命力/注意力/损伤）；\textbf{maintain} 强调维持现状/标准（\textit{maintain a road} 养路；\textit{maintain order} 维持秩序）；\textbf{support} 强调支撑/支持（物理/经济/情感）；\textbf{uphold} 强调维护原则/正义（\textit{uphold the law}）。在"维持生计"义上 sustain 与 support 部分互通。
\end{gram}

\begin{origin}
13 世纪借自古法语 \textit{sostenir}，拉丁语 \textit{sustinere}（支撑、维持）：\textit{sus-}（sub- 的同化形式"在下"）+ \textit{tenere}（握、持）。同源词 \textbf{sustainable}（可持续的）、\textbf{sustenance}（食物、生计）、\textbf{sustained}（持续的）、\textbf{tenable}（站得住脚的）、\textbf{abstain}（ab- 离开 + tenere $\to$ 戒绝）、\textbf{detain}（拘留）。
\end{origin}

\vw{manner}{/ˈmænə(r)/}{n.}

\sense{n.}{方式，方法}
\col{in a ... manner}
\ex{Please reply in a polite manner — the customer is easily offended.}
\col{in the manner of sb/sth 以……风格}
\ex{The poem is written in the manner of the great Tang poets.}

\sense{n.}{举止}
\col{have a ... manner}
\ex{She has a warm, open manner that immediately puts people at ease.}

\sense{n.}{(manners) 礼貌}
\col{It is good/bad manners to do sth}
\ex{It is bad manners to talk with your mouth full.}
\col{have no manners 没礼貌}
\ex{Their children have no manners at all — they've never been taught to say please and thank you.}

\begin{gram}{manner(s) 的两种"礼貌"用法}
\textbf{singular manner}（举止、态度）：\textit{a friendly manner}（友好的态度）、\textit{to the manner born}（生来就习惯/适于）。\textbf{plural manners}（礼貌、礼仪）：\textit{It is good manners to...}（做……是有礼貌的）、\textit{mind your manners}（注意礼貌）、\textit{comic/table manners}（喜剧/餐桌礼仪）。莎士比亚《哈姆雷特》"to the manner born"现已成英语成语。
\end{gram}

\begin{origin}
古法语 \textit{maniere}（方式、习惯），拉丁语 \textit{manuarius}（属于手的），源于 \textit{manus}（手）。同源词 \textbf{manual}（手工的；手册）、\textbf{manufacture}（制造；手做之物）、\textbf{manoeuvre}（机动；策略；演习）、\textbf{manage}（管理；驾驭）、\textbf{maintain}（用手撑住$\to$维持）、\textbf{mannerism}（习气、特殊癖性）。"方式"出自"手的动作"，是工具理性的具象化。
\end{origin}

\vw{consumption}{/kənˈsʌmpʃn/}{n.}

\sense{n.}{消耗，消耗量}
\col{fuel/water/energy consumption}
\ex{The new car model has a much lower fuel consumption than its predecessor.}

\sense{n.}{消费}
\col{domestic/personal consumption 国内/个人消费}
\ex{Rising domestic consumption has helped drive the country's economic recovery.}
\col{mass consumption 大众消费}
\ex{The industrial revolution gave rise to a new era of mass consumption.}

\begin{origin}
15 世纪借自拉丁语 \textit{consumptionem}（使用、消耗），源于 \textit{consumere}（用尽、消耗）：\textit{con-}（完全）+ \textit{sumere}（拿、取）。同源词 \textbf{consume}（消耗；吃光；烧尽）、\textbf{consumer}（消费者）、\textbf{consumerism}（消费主义）、\textbf{assume}（承担；假定）、\textbf{resume}（恢复；摘要；\textit{re-} 再 + \textit{sumere} $\to$再拿起）、\textbf{presume}（推测；\textit{prae-} 前 + \textit{sumere} $\to$预先拿$\to$假设）。
\end{origin}

\vw{church}{/tʃɜːtʃ/}{n.}

\sense{n.}{教堂}
\col{at/in/to church}
\ex{They go to church every Sunday morning.}

\sense{n.}{礼拜}
\ex{Most of the villagers were still at church when the storm hit.}

\begin{gram}{church 与 the Church 的"机构化"用法}
无冠词小写 \textbf{church}（教堂；做礼拜）；大写加 the \textbf{the Church}（教会这一组织）——既可指整个基督教教会（\textit{the Catholic Church}），也可指特定宗派（\textit{the Anglican Church}）。这是英语宗教词汇"小写=建筑/活动，大写=机构"的标准对应。
\end{gram}

\begin{origin}
古英语 \textit{circe}（教堂），希腊语 \textit{kyriakon}（主的），源于 \textit{kyrios}（主）。同源词 \textbf{kirk}（苏格兰英语"教堂"）、德语 \textit{Kirche}（教堂）、\textbf{cherub}（基路伯；天使，词源同 kyrios）。希腊语原义"属于主的房屋"，日耳曼语系同源词族非常整齐。
\end{origin}

\vp{come down to}

\sense{短语}{可归结为，可归纳为}
\col{come down to sth}
\ex{When you strip away all the legal jargon, the dispute comes down to a question of ownership.}
\ex{It all comes down to whether we can afford the extra time.}

\begin{gram}{come down to + 名词/动名词：表"实质是"的口语化逻辑归结}
\textbf{come down to} 后接名词/动名词短语，结构等效于 "essentially be a matter of..."。同源表达：\textit{What it boils down to is...}（煮到最后$\to$归结为）、\textit{The bottom line is...}（底线是）。三者皆口语"实质化"用语，正式写作宜用 \textit{essentially} / \textit{in essence}。
\end{gram}

\vw{code}{/kəʊd/}{n.}

\sense{n.}{行为规范，道德准则}
\col{a code of conduct/behaviour/ethics}
\ex{All soldiers are expected to follow a strict code of conduct.}
\col{a moral code 道德准则}
\ex{The villagers live by a strict moral code passed down from generation to generation.}

\sense{n.}{法规，法典}
\col{criminal/civil/penal code}
\ex{The new penal code came into force at the beginning of the year.}

\sense{n.}{密码，代码}
\col{break/crack a code 破译密码}
\ex{The Allies worked day and night to break the German Enigma code.}
\col{barcode/area code/postcode}
\ex{Enter the postcode to find the nearest store.}

\begin{gram}{code 的复合词族（着装规定）}
\textbf{dress code}（着装规定：\textit{a strict dress code}）、\textbf{zip code}（邮编，美式）、\textbf{morse code}（莫尔斯电码）、\textbf{genetic code}（遗传密码）。dress code 常见于酒店/餐厅/夜店告示：\textit{smart-casual dress code}（半正式休闲）、\textit{black-tie dress code}（黑领结正装）。
\end{gram}

\begin{origin}
13 世纪借自古法语 \textit{code}，拉丁语 \textit{codex}（树干；木板书；法典）。同源词 \textbf{codify}（编纂；整理成法典）、\textbf{codicil}（遗嘱附录）、\textbf{codex}（古籍抄本）。"密码"义 19 世纪起自电报而生——莫尔斯码将字母转化为点划，与"将文字编码为法律条款"是同一逻辑。
\end{origin}

\vw{item}{/ˈaɪtəm/}{n.}

\sense{n.}{一件商品（或物品）}
\col{item of clothing/furniture/equipment}
\ex{She packed a few essential items of clothing into a small bag.}

\sense{n.}{项目，条款}
\col{item on the agenda/list}
\ex{Let's move on to the next item on the agenda.}
\col{item by item 逐条}
\ex{The contract was reviewed item by item by the legal team.}

\begin{gram}{item 作为副词的"新闻条目"用法}
古英语中 item 作副词"同样地、亦"（\textit{item, one must remember that...}），后衍生"分条列出"之义——这是中世纪清单"第一条……第二条……"的句首副词。今已不通用，但现代英语"item"一词作为"分条列出之条目"的本义即源于此。
\end{gram}

\begin{origin}
14 世纪借自拉丁语 \textit{item}（同样地、亦；列举时"又、其次"），最终可能与 \textit{ita}（如此）相关。同源词 \textbf{iterate}（迭代；重申）。item 从"亦、再"演变为"所列之项"，是中世纪账簿与法律文书的句法化石。
\end{origin}

\vw{elegant}{/ˈelɪɡənt/}{adj.}

\sense{adj.}{漂亮雅致的，精美的}
\col{an elegant dress/dress/hotel}
\ex{She wore an elegant black dress to the gala dinner.}
\col{an elegant solution/solution 巧妙的解决方案}
\ex{The mathematician found an elegant solution to a problem that had puzzled experts for decades.}

\sense{adj.}{优美的，文雅的}
\col{an elegant speaker/style/manner}
\ex{His essays are famous for their elegant prose and subtle wit.}

\begin{lex}{elegant / graceful / refined / sophisticated}
\textbf{elegant} 强调简洁精致、线条/设计优美（多用于衣着/方案）；\textbf{graceful} 强调动作/姿态自然柔美（多用于人体/动物）；\textbf{refined} 强调经打磨变得精致（多用于品味/举止）；\textbf{sophisticated} 强调世故、复杂老练（多用于人/技术）。an elegant solution 不可说成 a graceful solution——优雅方案而非优雅动作。
\end{lex}

\begin{origin}
15 世纪借自拉丁语 \textit{elegans}（有鉴别力的、雅致的），是 \textit{eligere}（挑选）的现在分词：\textit{ex-}（出）+ \textit{legere}（选、拣）。同源词 \textbf{elect}（选举；选出）、\textbf{election}（选举）、\textbf{eligible}（合格的；可被选的）、\textbf{intelligent}（\textit{inter-} 间 + legere $\to$能于其间挑选$\to$聪明的）、\textbf{lecture}（演讲；\textit{legere} 读）、\textbf{legend}（传奇；\textit{legere} 读$\to$应读之事）。"雅致"源于"会挑选"。
\end{origin}

\vp{make one's way to}

\sense{短语}{去，前往}
\col{make one's way to ...}
\ex{After the lecture, the students made their way to the dining hall.}
\ex{She made her way through the crowded market to the riverbank.}

\begin{gram}{way 在运动短语中的意象}
英语表"行进"的短语多以 way 为核心：\textbf{make one's way}（前往）、\textbf{find one's way}（找到路）、\textbf{lose one's way}（迷路）、\textbf{push one's way}（挤过去）、\textbf{feel one's way}（摸索前进）、\textbf{wend one's way}（蜿蜒前进）、\textbf{take one's way}（出发）。皆"自己 + way"的结构，way 充当"行进路径"的容器。
\end{gram}

\vw{souvenir}{/ˌsuːvəˈnɪə(r)/}{n.}

\sense{n.}{纪念物，纪念品}
\col{a souvenir of/from sth}
\ex{I bought a small souvenir of my trip to Edinburgh — a miniature castle made of local stone.}
\col{a souvenir shop}
\ex{The narrow streets are lined with souvenir shops selling postcards and trinkets.}

\begin{origin}
18 世纪借自法语 \textit{souvenir}（记忆、纪念品），是动词 \textit{souvenir}（记起）的名词形式，晚期拉丁语 \textit{subvenire}（来到心中、记起）：\textit{sub-}（在下面）+ \textit{venire}（来）。同源词 \textbf{subvention}（补助金——"从下面来支持"）、\textbf{intervene}（干预；inter- 中间 + venire）。"纪念品"即"某物来到心中"之物。
\end{origin}

\vw{antique}{/ænˈtiːk/}{n. \& adj.}

\sense{n.}{文物，古董}
\col{an antique dealer/shop}
\ex{My grandmother loves browsing the local antique shop on Saturday mornings.}

\sense{adj.}{古老的，古董的}
\col{antique furniture/jewellery}
\ex{The hotel is full of antique furniture from the Victorian era.}
\col{in antique style 古色古香}
\ex{The little chapel is decorated in antique style with hand-carved wooden screens.}

\begin{gram}{antique / ancient / antiquated 三者年龄感}
\textbf{antique} 通常指 100 年以上、可作收藏的器物（\textit{an antique clock} 古董钟）；\textbf{ancient} 形容极古老的文明/时期（\textit{ancient Egypt} 古埃及）；\textbf{antiquated} 是贬义形容词，形容事物"过时的、老掉牙的"（\textit{antiquated ideas} 落伍观念）。ancient 偏中性，antiquated 偏贬义。
\end{gram}

\begin{origin}
16 世纪借自拉丁语 \textit{antiquus}（古老的、古代的），与 \textit{ante}（在前）相关。同源词 \textbf{antiquity}（古代；古迹）、\textbf{antiquarian}（研究古物的；古物研究者）、\textbf{antiquated}（过时的）。注意 antique（形容词）重音在前，antique（名词）重音亦在前，但作"古风/古色古香"时两种重音都常见。
\end{origin}

\vw{dune}{/djuːn/}{n.}

\sense{n.}{沙丘}
\col{a sand/sandy dune}
\ex{The wind keeps shifting the sand dunes along the coast.}
\col{dune buggy 沙滩越野车}
\ex{They spent a wild afternoon racing dune buggies over the soft sand.}

\begin{origin}
18 世纪借自法语 \textit{dune}（沙丘），中古荷兰语 \textit{dūne}（沙丘；今荷兰语 \textit{duin}），原始日耳曼语 *\textit{dunon}（高地、丘陵）。同源词 \textbf{down}（丘陵——古英语 \textit{dūn}，与 dune 同源；\textit{the South Downs} 英国南唐斯丘陵）、\textbf{duneland}（沙丘地带）、 \textbf{ dune grass}（沙丘草）。荷兰低地之国，故沙丘词借自荷兰语。
\end{origin}

\vw{fragile}{/ˈfrædʒaɪl/}{adj.}

\sense{adj.}{脆弱的}
\col{a fragile economy/alliance/peace}
\ex{The peace agreement is still fragile and could collapse at any moment.}
\col{feel fragile 感觉虚弱}
\ex{She felt extremely fragile after the long flight.}

\sense{adj.}{易碎的}
\col{fragile glassware/porcelain}
\ex{Please handle the boxes with care — they contain fragile items.}

\begin{gram}{fragile / brittle / delicate / frail}
\textbf{fragile} 强调易碎、脆弱（物品/局势/身体）；\textbf{brittle} 强调硬而易断裂（多指金属/糖/玻璃，受力即裂）；\textbf{delicate} 强调精致、需小心对待（\textit{a delicate flower} 娇嫩的花；\textit{a delicate balance} 精妙的平衡）；\textbf{frail} 强调瘦弱/不强壮（人多指身体、声音）。四者中文皆可译"脆弱"，但适用对象不同。
\end{gram}

\begin{origin}
16 世纪借自拉丁语 \textit{fragilis}（易碎的），源于 \textit{frangere}（破、折）。同源词 \textbf{fragment}（碎片）、\textbf{frangible}（可破碎的）、\textbf{fraction}（分数；小部分）、\textbf{fracture}（骨折）、\textbf{refrain}（节制；\textit{re-} 回 + frangere $\to$ 折回$\to$忍住）、\textbf{infringe}（侵犯；\textit{in-} 入 + frangere $\to$ 折入他人领地）。
\end{origin}

\vw{resolution}{/ˌrezəˈluːʃn/}{n.}

\sense{n.}{决心，决定}
\col{make/keep/break a resolution}
\ex{I made a New Year's resolution to read at least one book a month.}
\col{a man/woman of great resolution 意志坚定者}
\ex{She approached the task with calm and great resolution.}

\sense{n.}{决议}
\col{pass/adopt a resolution}
\ex{The UN General Assembly passed a resolution calling for an immediate ceasefire.}
\col{a UN Security Council resolution}
\ex{The Security Council resolution demanded the withdrawal of all foreign troops.}

\sense{n.}{解决}
\col{the resolution of a problem/conflict/dispute}
\ex{We are hoping for a peaceful resolution of the conflict.}
\col{conflict resolution 冲突解决}
\ex{The two sides have agreed to seek international help in the resolution of their dispute.}

\begin{gram}{resolution 三义中文皆不同：决心 / 决议 / 解决}
\textbf{resolution} 是英语"一词三义"经典案例：1. 决心（个人：\textit{New Year resolution}）；2. 决议（机构：\textit{UN resolution}）；3. 解决（问题：\textit{conflict resolution}）。三义皆源于拉丁语 \textit{resolvere}（解开、解决）——个人"下定决心"是"解开心中疑团"，机构"通过决议"是"解开议程难点"，社会"解决冲突"是"解开纠葛"。
\end{gram}

\begin{origin}
15 世纪借自拉丁语 \textit{resolutionem}（解开、解决），源于 \textit{resolvere}：\textit{re-}（回）+ \textit{solvere}（松开、解开）。同源词 \textbf{resolve}（决心；解决；分解）、\textbf{resolved}（下定决心的）、\textbf{soluble}（可溶的；可解决的）、\textbf{absolute}（绝对的；\textit{ab-} 离开 + solvere $\to$ 完全松开$\to$不受约束的）、\textbf{dissolve}（溶解；解散）、\textbf{solve}（解决；解开）。
\end{origin}

\vw{grain}{/ɡreɪn/}{n.}

\sense{n.}{细粒，颗粒}
\col{a grain of sand/salt/sugar}
\ex{There wasn't a grain of truth in what he said.}
\col{against the grain 逆纹理；违反本性}
\ex{It goes against the grain for him to admit he was wrong.}

\sense{n.}{谷物}
\col{grain harvest/storage}
\ex{The country is one of the world's largest grain exporters.}
\col{whole grains 全谷物}
\ex{Health experts recommend eating more whole grains and less processed food.}

\begin{gram}{grain of truth/sense/hope 等"小剂量"搭配}
英语中 a grain of 后接抽象名词，传达"极少、微量"之意：\textit{a grain of truth}（一丝真实）、\textit{a grain of comfort}（一丝安慰）、\textit{a grain of hope}（一线希望）、\textit{a grain of sense}（一点点道理）。同构搭配：\textit{a grain of salt}（保留态度，"带点盐"$\to$不轻信）。
\end{gram}

\begin{origin}
13 世纪借自古法语 \textit{grain}，拉丁语 \textit{granum}（谷粒、种子），原始印欧语 *\textit{gerə-}（成熟、变老、生长）。同源词 \textbf{grainy}（有颗粒感的）、\textbf{granary}（粮仓）、\textbf{granule}（小颗粒）、\textbf{grind}（磨）、\textbf{engrain}（使根深蒂固）、\textbf{corn}（谷物——印欧 *\textit{gerə-} 与日耳曼 *\textit{kurnan} 同源）。"谷物"与"生长"在印欧语里同根——种子是"已成熟的果实"。
\end{origin}

\vp{beat down}

\sense{短语}{（阳光）强烈照射，曝晒}
\ex{The afternoon sun beat down on the desert, making the sand too hot to walk on.}
\ex{We sheltered under a tree while the sun beat down mercilessly.}

\begin{gram}{beat 短语小家族}
\textbf{beat down}（晒；压价）；\textbf{beat up}（痛打）；\textbf{beat off}（击退）；\textbf{beat about the bush}（拐弯抹角）；\textbf{beat a retreat}（败退；匆匆撤退）；\textbf{off the beaten track}（人迹罕至）；\textbf{heart beating fast}（心跳加速）。beaten 形式：\textit{beaten record}（被打破的纪录）、\textit{beaten path}（踏出的路）。
\end{gram}

\vw{contrast}{/kənˈtrɑːst/}{vi. \& vt. \& n.}

\sense{vi.}{形成对比}
\col{contrast with sth}
\ex{His confident manner contrasts sharply with his brother's shyness.}
\ex{The bright red door contrasts beautifully with the white walls of the house.}

\sense{vt.}{对比，对照}
\col{contrast A with/against B}
\ex{The essay contrasts the poet's early work with his later, more political writings.}

\sense{n.}{差异，差别}
\col{in contrast to/with sth 与……形成对比}
\ex{In contrast to her outgoing sister, Mary is quiet and reserved.}
\col{a contrast between A and B}
\ex{There is a striking contrast between the two candidates' views on the economy.}

\sense{n.}{对照物}
\col{be a contrast to sth}
\ex{The new building is a stark contrast to the old houses around it.}

\begin{gram}{contrast 重音转移：动词 vs. 名词}
\textbf{contrast} 作动词重音在后 /kənˈtrɑːst/（"对比"），作名词重音在前 /ˈkɒntrɑːst/（"差异、对照"）。同类的重音转移词：\textbf{record}（v. /rɪˈkɔːd/ 录制；n. /ˈrekɔːd/ 记录）、\textbf{conduct}（v. /kənˈdʌkt/ 实施；n. /ˈkɒndʌkt/ 行为）、\textbf{increase}（v. /ɪnˈkriːs/ 增加；n. /ˈɪŋkriːs/ 增长）。
\end{gram}

\begin{origin}
15 世纪借自法语 \textit{contraster}（相对而立），晚期拉丁语 \textit{contrastare}（相对、抵抗）：\textit{contra-}（反对、相对）+ \textit{stare}（站立）。同源词 \textbf{contrastive}（对比的）、\textbf{contrastive analysis}（对比分析）、\textbf{contra}（反对）、\textbf{contrary}（相反的）、\textbf{encounter}（遭遇；\textit{en-} 入 + \textit{contra} $\to$ 迎面相遇）、\textbf{constant}（恒定的；\textit{con-} 一起 + \textit{stare} $\to$一起站着$\to$稳固）。"对比"源自"两物相对而立"。
\end{origin}

\vw{ray}{/reɪ/}{n.}

\sense{n.}{光线}
\col{a ray of sunlight/light}
\ex{The first rays of sunlight crept through the curtains.}
\col{a ray of hope 一线希望}
\ex{After months of uncertainty, there is finally a ray of hope for the hostages.}

\sense{n.}{射线}
\col{gamma/ultraviolet/cosmic rays}
\ex{Cosmic rays constantly bombard the earth's atmosphere from outer space.}

\begin{lex}{ray / beam / shaft 的"光"细分}
\textbf{ray} 强调"从一点射出的细光线"（多用于阳光、X 射线、宇宙射线）；\textbf{beam} 强调"较宽的光柱"（\textit{a beam of light} 探照灯光柱、汽车车灯）；\textbf{shaft} 强调"穿过缝隙的光束"（\textit{a shaft of sunlight} 穿过缝隙的阳光）。三个词皆有"光"的具象性，但光的形态不同。
\end{lex}

\begin{origin}
14 世纪借自古法语 \textit{rai}，拉丁语 \textit{radius}（光线、辐条、半径）。同源词 \textbf{radius}（半径）、\textbf{radiate}（辐射；发射）、\textbf{radiator}（暖气片；散热器）、\textbf{radial}（辐射状的）、\textbf{radio}（无线电；通过辐射传播）、\textbf{radium}（镭——放射性元素）、\textbf{rayon}（人造丝）。"光线"与"半径"在拉丁语中同根——半径就是从圆心射出的"光线"。
\end{origin}

\vw{chew}{/tʃuː/}{vi. \& vt.}

\sense{vi. \& vt.}{咀嚼，嚼碎}
\col{chew (on) sth}
\ex{Chew your food well before swallowing.}
\col{chewing gum 口香糖}
\ex{A sign on the wall read: ``No chewing gum allowed.''}

\sense{vi. \& vt.}{咬住}
\col{chew on sth 反复咀嚼（亦引申为反复思考）}
\ex{The puppy chewed on the slipper until it was ruined.}
\ex{I'll need a few days to chew on your proposal before giving you an answer.}

\begin{gram}{chew over / chew on：反复思考的隐喻}
\textbf{chew on/over}（反复思考、仔细琢磨）、\textbf{chew the fat}（闲聊、唠嗑）、\textbf{bite off more than one can chew}（贪多嚼不烂）。chew 在英语中既指嘴巴的物理动作，也隐喻为大脑"反复咀嚼"信息。
\end{gram}

\begin{origin}
古英语 \textit{ceowan}（咀嚼），原始日耳曼语 *\textit{kewan}。同源词 \textbf{chewy}（耐嚼的）、古英语 \textit{cēow}（咀嚼）。jaw（颚）与 chew 同源——jaw（颚）即"咬东西的部位"。
\end{origin}

\vw{starve}{/stɑːv/}{vi. \& vt.}

\sense{vi. \& vt.}{(使) 饿死，(使) 挨饿}
\col{starve to death}
\ex{Hundreds of refugees starved to death during the long siege.}
\col{be starving 饿得要死}
\ex{I'm absolutely starving — when is dinner ready?}

\sense{vi. \& vt.}{(使) 严重缺乏}
\col{be starved of sth}
\ex{The remote village is starved of basic medical care and clean water.}
\ex{The arts have been starved of funding for years.}

\begin{gram}{starve of / be starved of：被剥夺的"饥饿"隐喻}
\textbf{starve sb of sth}（使某人缺乏/渴望某物）：\textit{The children were starved of affection}（孩子们渴望关爱）。与 \textbf{be deprived of} 同义但更具情感色彩。同构：\textit{a garden starved of water}（严重缺水的花园）、\textit{a mind starved of stimulation}（缺乏刺激的大脑）。
\end{gram}

\begin{origin}
古英语 \textit{steorfan}（死、僵死），原始日耳曼语 *\textit{sterbanan}（僵死）。同源词 \textbf{starveling}（挨饿者；饿瘦的人或动物）、古英语 \textit{steorfan}。Starve 原义"死"——饿死是"死"的最常见形式；现代英语"饿死"义成主流，古义"僵死"已不存。
\end{origin}

\vw{departure}{/dɪˈpɑːtʃə(r)/}{n.}

\sense{n.}{离开，起程}
\col{arrival and departure 抵离}
\ex{Please check the departure board for the latest gate information.}
\col{a departure from ... 偏离、背离}
\ex{His latest novel marks a sharp departure from his earlier, more realistic style.}

\begin{gram}{departure / arrival：动词对应的名词化}
\textbf{depart}（v. 离开）$\to$ \textbf{departure}（n.）；\textbf{arrive}（v. 到达）$\to$ \textbf{arrival}（n.）。英语 -ure 后缀常将动词名词化（\textit{depart/ure, fail/ure, mix/ure, pres/ure}），与 -tion 互补（\textit{admit/admission}）。注意：\textit{depart} 是不及物动词，不可说 \textit{depart the room}，应说 \textit{depart from the room}；但 \textit{departure from} 即由这一结构名词化而来。
\end{gram}

\begin{origin}
15 世纪借自古法语 \textit{departir}（分开、离开），拉丁语 \textit{dispertire}（分开）：\textit{dis-}（分开）+ \textit{partire}（分、分割）。同源词 \textbf{depart}（离开）、\textbf{department}（部门；分开的一部分）、\textbf{particle}（微粒）、\textbf{partner}（伙伴——分摊的"那一份"）、\textbf{particular}（特别的；"分成小份"$\to$详尽的）、\textbf{part}（部分；分开）、\textbf{partition}（分隔；瓜分）。
\end{origin}

\vw{border}{/ˈbɔːdə(r)/}{n. \& vi. \& vt.}

\sense{n.}{国界，边界}
\col{cross the border}
\ex{They crossed the border into France at dawn.}
\col{a border town/area 边境小镇/地区}
\ex{The border area has been heavily militarised in recent months.}

\sense{vi. \& vt.}{(与……) 接壤}
\col{border (on/upon) sth}
\ex{My country borders on three different countries.}
\ex{The park borders the river for several kilometres.}

\begin{gram}{border / boundary / frontier / rim 的"边界"细分}
\textbf{border} 强调政治/地理边界线（\textit{cross the border}）；\textbf{boundary} 强调抽象/制度界限（\textit{the boundary between right and wrong}）；\textbf{frontier} 强调"边疆"——未开发的边远地带或国境的边缘地区（\textit{the American frontier} 美国西部边疆）；\textbf{rim} 强调圆形物的边缘（\textit{the rim of a cup} 杯口）。四者中文皆可译"边界"，但侧重点不同。
\end{gram}

\begin{origin}
14 世纪借自古法语 \textit{bord}（边、岸），最终可能溯至古日耳曼语 *\textit{bord}（板、边）。同源词 \textbf{board}（板；登上）、\textbf{embark}（登船）、\textbf{debark}（下船）、\textbf{disembark}（上岸）、\textbf{overboard}（越过船舷）。"边界"原义是"木板、岸边"——边境上的界碑与水边之岸，皆是"边"之物化。
\end{origin}

\vw{hut}{/hʌt/}{n.}

\sense{n.}{(简陋的) 小屋，棚屋}
\col{a mountain/shepherd's/beach hut}
\ex{We took shelter from the storm in a tiny mountain hut.}
\col{a mud hut 泥屋}
\ex{The villagers still live in traditional mud huts with thatched roofs.}

\begin{origin}
17 世纪借自法语 \textit{hutte}（小屋），中世纪高地德语 \textit{hutte}（小屋、窝棚）。同源词德语 \textit{Hütte}（小屋）、荷兰语 \textit{hut}（棚屋）。与"帽子"（hat）无词源关系——hut 的核心是"棚"，hat 的核心是"罩"。
\end{origin}

\vw{dot}{/dɒt/}{vt. \& n.}

\sense{vt.}{遍布；使布满，点缀}
\col{be dotted with sth}
\ex{The hillside was dotted with wild flowers in every colour of the rainbow.}
\ex{The lake is dotted with small fishing boats.}

\sense{vt.}{加点}
\col{dot one's i's 在 i 上加点（亦指一丝不苟）}
\ex{Please dot your i's so the file is legible.}

\sense{n.}{点}
\col{a dot on the page/horizon}
\ex{The boat was now just a tiny dot on the horizon.}
\col{on the dot 准时}
\ex{She arrived at the meeting at three o'clock on the dot.}

\begin{gram}{on the dot / in a row / in no time：副词性小品的强化语}
英语口语中许多小品词起强调作用：\textbf{on the dot}（正点；准时）、\textbf{year after year}（年复一年）、\textbf{in no time}（立刻）、\textbf{off and on}（断断续续）、\textbf{over and over}（反复）。这些短语中 dot/row/time/on/over 都不取本义，仅作为语法骨架。
\end{gram}

\begin{origin}
古英语 \textit{dott}（头顶的肿块；一点），可能与古英语 \textit{dyttan}（堵住）相关。同源词 \textbf{dottle}（烟斗中残留的烟丝——堵在管底）、\textbf{dotty}（点多的；古怪的）。注意 \textit{dot the i's and cross the t's}（一丝不苟）——这两个"点笔画"是英文书法的基础。
\end{origin}

\vw{flap}{/flæp/}{vi. \& vt.}

\sense{vi. \& vt.}{拍打，摆动}
\col{flap in the wind}
\ex{The sails flapped gently in the breeze.}
\ex{The tent flap flapped noisily in the strong wind.}

\sense{vi. \& vt.}{振（翅）}
\col{flap its wings}
\ex{The heron flapped its wings slowly and rose from the marsh.}

\begin{origin}
14 世纪拟声词，模拟拍打声。同源词 \textbf{flapper}（拍打者；20 年代新女性）、\textbf{floppy}（松垮的）、\textbf{flop}（扑通落下；失败）。flap、flop、flip、flutter、flatter 一族皆为拟声：模拟物体拍打、翻转、轻拍之声。英语拟声词族是研究语言音响象征的富矿。
\end{origin}

\vw{bunch}{/bʌntʃ/}{n.}

\sense{n.}{串，束}
\col{a bunch of flowers/grapes/bananas}
\ex{She bought a bunch of red roses for her mother's birthday.}
\col{a bunch of keys 一串钥匙}
\ex{He jingled a bunch of keys as he walked towards the door.}

\sense{n.}{大量，许多（口语）}
\col{a bunch of sth}
\ex{There were a whole bunch of new students in the class this term.}
\col{the best of a bad bunch 一群不太行中的最佳}
\ex{Among the candidates, he was the best of a bad bunch.}

\begin{gram}{a bunch of 后接可数名词复数时动词用单/复数皆可}
正式写作：\textit{a bunch of students are waiting}（一群学生在等）——以"学生"为焦点，复数谓语。口语与非正式：\textit{a bunch of students is waiting}（一群学生在等）——以"群"为焦点，单数谓语。同类的"集合名词 + of + 复数"：\textit{a group of, a team of, a flock of}。
\end{gram}

\begin{origin}
14 世纪借自古法语 \textit{bunch}（肿块、瘤），可能与古法语 \textit{bombe}（肿）相关。同源词 \textbf{bundle}（捆、束；与 bunch 形近但来自古英语 \textit{byndelle}）。bunch 与 bundle 常可互换，但 bunch 偏指自然成串之物（葡萄、香蕉、钥匙），bundle 偏指人工捆扎之物（木柴、文件、衣物）。
\end{origin}

\vw{fetch}{/fetʃ/}{vt.}

\sense{vt.}{(去) 拿来，(去) 请来}
\col{fetch sth (for sb)}
\ex{Could you fetch me a glass of water from the kitchen?}
\ex{She fetched the doctor immediately when the child fell ill.}

\sense{vt.}{卖得（多少钱）}
\col{fetch a price}
\ex{The painting fetched a record price at the auction.}

\begin{gram}{fetch / bring / take / carry 的"取送"细分}
\textbf{fetch}（去取然后拿来——包含去和回两个动作）、\textbf{bring}（带来——从别处向说话人方向）、\textbf{take}（带走——从说话人方向向别处）、\textbf{carry}（搬运——无方向偏好，可远可近）。最易错的是 fetch vs. bring：\textit{Fetch me a chair from the next room.}（去隔壁拿把椅子来）。
\end{gram}

\begin{origin}
古英语 \textit{fecan}（取来、带来），是 \textit{fetian}（取来）的南方方言变体。同源词 \textbf{far-fetched}（牵强的；远道取来的）、德语 \textit{fassen}（抓住）。现代英语 fetch 是"取+回"复合动作，是"单向"动词中最具动作感的一个。
\end{origin}

\vw{possession}{/pəˈzeʃn/}{n.}

\sense{n.}{私人物品，个人财产}
\col{possessions（复数用法）}
\ex{She arrived at the hotel with only a few personal possessions in a small bag.}
\col{a man/woman of few possessions 财产不多的人}
\ex{Despite his fame, he remained a man of few possessions.}

\sense{n.}{具有，拥有}
\col{be in possession of sth 拥有某物}
\ex{The defendant is currently in possession of the disputed property.}
\col{come into possession of sth 获得某物}
\ex{She came into possession of a small fortune from a distant relative.}
\col{take possession of sth 占有某物}
\ex{The new owners took possession of the house at the beginning of the month.}

\begin{gram}{in possession of / in the possession of：主动 vs. 被动}
\textbf{in possession of sth}（某人拥有某物，主动含义）：\textit{He is in possession of the documents.}（他持有这些文件）。\textbf{in the possession of sb}（某物被某人拥有，被动含义）：\textit{The documents are in the possession of the police.}（这些文件被警方持有）。中间加 the 是英语"所有关系"经典标记。
\end{gram}

\begin{origin}
14 世纪借自古法语 \textit{possession}，拉丁语 \textit{possessionem}（占据、拥有），源于 \textit{possidere}（占据、拥有）：\textit{pot-}（主人；与 \textit{power} 同源）+ \textit{sedere}（坐）。同源词 \textbf{possess}（拥有）、\textbf{possessive}（所有的；占有欲强的）、\textbf{possible}（可能的；\textit{pot-} + \textit{esse} $\to$能做主的$\to$能做的）。"拥有"在拉丁词源中即"主人坐下"——所有权是"坐下来的姿态"。
\end{origin}

\vw{intently}{/ɪnˈtentli/}{adv.}

\sense{adv.}{热切地，专注地}
\col{stare/listen intently}
\ex{The child stared intently at the magician's hands, trying to spot the trick.}
\ex{The jury listened intently as the witness gave her account of the incident.}

\begin{lex}{intent / intention / intently / intentional 的派生族}
\textbf{intent}（n. 意图；adj. 专注的、热切的）、\textbf{intention}（n. 意图、打算）、\textbf{intently}（adv. 专注地）、\textbf{intentional}（adj. 故意的）。区别：\textit{intent on doing sth}（专心做某事；形容词）；\textit{intention to do sth}（做某事的打算；名词）。\textit{with intent to kill}（蓄意杀人）——法律用语。
\end{lex}

\begin{origin}
拉丁语 \textit{intentus}（伸展、专心、紧张），是 \textit{intendere}（伸向、指向、打算）的过去分词：\textit{in-}（向内）+ \textit{tendere}（伸、拉伸）。同源词 \textbf{intend}（打算）、\textbf{intense}（强烈的）、\textbf{intensive}（密集的）、\textbf{pretend}（假装；\textit{prae-} 在前 + tendere $\to$在前面伸出$\to$把意图置于行为前）、\textbf{tend}（倾向）、\textbf{attend}（注意；出席）。"专注"是"将注意力伸展到一点"——认知作为动作的隐喻。
\end{origin}

\vw{hike}{/haɪk/}{n. \& vi. \& vt.}

\sense{n.}{远足，徒步旅行}
\col{go on/for/take a hike 去远足}
\ex{We went on a long hike through the mountains at the weekend.}
\col{a five-mile/half-day hike}
\ex{There's a lovely five-mile hike along the coast from the village.}

\sense{vi. \& vt.}{去……远足，做徒步旅行}
\col{hike in the mountains/hills}
\ex{They hiked for hours through the dense forest.}
\col{hike a trail 沿小路徒步}
\ex{He has hiked the entire Appalachian Trail over the summer.}

\begin{gram}{hike 作动词的"涨价"用法（拓展义）}
\textbf{price hike}（涨价）、\textbf{tax hike}（加税）、\textbf{rent hike}（租金上涨）。hike 的本义是"长途跋涉"，引申为"价格/数量大幅上升"——隐喻基础是"向上跋涉"。
\end{gram}

\begin{origin}
19 世纪方言词，可能与 \textit{yike}（踱步、来回走）相关，或拟声于"嗨哟嗨哟"行走之声。hiker（远足者）、hiking boots（登山靴）皆同期出现。hike 作为"涨价"是 20 世纪美式英语的新引申。
\end{origin}

\vw{awesome}{/ˈɔːsəm/}{adj.}

\sense{adj.}{令人惊叹的，使人惊惧的}
\col{an awesome sight/experience/power}
\ex{The view from the top of the mountain was truly awesome.}
\ex{We stood in silence before the awesome power of the waterfall.}

\sense{adj.}{很好的（口语）}
\ex{We had an awesome time at the concert.}
\col{awesome! 太好了！}
\ex{``I got the job!'' ``Awesome!''}

\begin{gram}{awesome / awful / awesome 的语体变迁}
\textbf{awesome} 原义"令人敬畏的"（awe + some）；现代美式口语中已泛化为"极好的"（类似 cool、amazing）。\textbf{awful} 原义"令人敬畏的"（awe + ful），现代义"糟糕的"。二者词源同根（awe 敬畏），但现代褒贬相反：awesome = 太棒了，awful = 太糟了。
\end{gram}

\begin{origin}
16 世纪由 \textit{awe}（敬畏）+ \textit{-some}（充满……的）合成。awe 来自古英语 \textit{ege}（恐惧、敬畏），原始日耳曼语 *\textit{agiz}。同源词 \textbf{awe}（敬畏）、\textbf{awful}（糟糕的；古义"令人敬畏的"）、\textbf{awesome}（令人敬畏的；今多义"极好的"）、\textbf{awkward}（笨拙的；古义"方向错误的，令人尴尬的"——\textit{awe} 词根与"反向"结合的产物）。
\end{origin}

\vw{liquid}{/ˈlɪkwɪd/}{adj. \& n.}

\sense{adj.}{液体的，液态的}
\col{liquid assets 流动资产（财务术语）}
\ex{The company needs to keep enough liquid assets to cover its short-term debts.}
\col{liquid nitrogen 液氮}
\ex{Liquid nitrogen is used to freeze biological samples quickly.}

\sense{n.}{液体}
\col{a clear/colourless/pale liquid}
\ex{Drink plenty of clear liquids when you have a fever.}

\begin{lex}{liquid / fluid / aqueous 的"流"细分}
\textbf{liquid} 强调"液态"——是 gas、solid 的对应（\textit{water is a liquid}）；\textbf{fluid} 强调"可流动"——既可指液体也可指气体（\textit{fluid mechanics} 流体力学）；\textbf{aqueous} 严格限于化学/医学，指"含水的、似水的"（\textit{an aqueous solution} 水溶液）。三者在科学语境中精确分工。
\end{lex}

\begin{origin}
14 世纪借自拉丁语 \textit{liquidus}（流动的、液体的），源于 \textit{liquere}（流动、清晰）。同源词 \textbf{liquefy}（液化）、\textbf{liquor}（酒；溶液）、\textbf{liquidate}（清偿；清算；\textit{liquid + -ate} $\to$ 化为现金）、\textbf{liquidator}（清算人）、\textbf{liqueur}（利口酒）、\textbf{listless}（无精打采的；\textit{lust} 欲望 + \textit{less} $\to$ 无欲望——与 liquid 形近但实不同源）。"液"与"清澈"在拉丁语中同根。
\end{origin}

\vw{fox}{/fɒks/}{n.}

\sense{n.}{狐狸}
\col{a fox's den 狐穴}
\ex{The fox's den is usually at the bottom of a deep burrow.}

\sense{n.}{狡猾的人}
\col{an old/a sly fox 老奸巨猾的人}
\ex{Watch out for him — he's a sly old fox.}

\begin{lex}{fox / vixen / fox cub 的"狐"细分}
\textbf{fox} 通用；\textbf{vixen} 母狐（亦可作贬义"泼妇"）；\textbf{fox cub / kit} 幼狐；雄狐在古英语有专门的 \textit{dog-fox}，今已不用。狐狸的"狡猾"形象在英语谚语中比比皆是：\textit{When the fox preaches, look to your geese.}（狐狸说教时看好你的鹅）。
\end{lex}

\begin{origin}
古英语 \textit{fox}，原始日耳曼语 *\textit{fuhsaz}（多毛的尾巴）。同源词德语 \textit{Fuchs}（狐狸）、荷兰语 \textit{vos}（狐狸）。fox 与拉丁语 \textit{vulpes}（狐，词族 \textit{vulpine} 狐狸似的、狡猾的）不同源——日耳曼语族与拉丁语族各有自己的"狐"字。
\end{origin}

\vw{polar}{/ˈpəʊlə(r)/}{adj.}

\sense{adj.}{极地的}
\col{the polar regions/ice caps}
\ex{The polar regions are warming at twice the global average rate.}
\col{a polar bear/expedition}
\ex{Polar bears are superbly adapted to life on the Arctic ice.}

\sense{adj.}{完全相反的}
\col{polar opposites}
\ex{The two candidates represent polar opposites on the political spectrum.}

\begin{origin}
16 世纪借自法语 \textit{polaire}，晚期拉丁语 \textit{polaris}（属于极的），源于拉丁语 \textit{polus}（极、地极），希腊语 \textit{ polos}（轴、极）。同源词 \textbf{pole}（极；杆）、\textbf{polarity}（极性）、\textbf{polarize}（两极分化）、\textbf{Polaroid}（偏光片；即时成像）、\textbf{poles apart}（天差地远）。"极"与"轴"在希腊语中同根——地轴的两端即"极"。
\end{origin}

\vw{delicacy}{/ˈdelɪkəsi/}{n.}

\sense{n.}{佳肴}
\col{a local/regional delicacy}
\ex{Fugu, the Japanese pufferfish, is considered a great delicacy — and a deadly one if prepared incorrectly.}

\sense{n.}{谨慎；细致}
\col{with delicacy}
\ex{The two leaders handled the sensitive issue with great delicacy.}
\col{a matter of delicacy 棘手的事情}
\ex{The question of pay is a matter of some delicacy in this household.}

\sense{n.}{敏感}
\col{the delicacy of a situation/instrument}
\ex{The delicacy of the political situation required careful handling.}

\begin{gram}{delicacy / delicate / delicious 的派生分叉}
\textbf{delicacy}（n. 佳肴；精妙；敏感）、\textbf{delicate}（adj. 精致的；娇嫩的；精妙的；微妙的）、\textbf{delicious}（adj. 美味的；愉快的）。注意：delicacy 表"佳肴"，delicious 表"美味"——delicacy 强调"需要小心对待的精致食物"，delicious 强调"感官愉悦"。
\end{gram}

\begin{origin}
14 世纪借自古法语 \textit{delicatesse}（精美、精致），源于拉丁语 \textit{delicatus}（使人愉悦的、娇嫩的），是 \textit{delectare}（使人愉悦）的过去分词。同源词 \textbf{delicate}（精致的；娇嫩的）、\textbf{delectable}（令人愉快的）、\textbf{delight}（高兴；\textit{de-} 加强 + \textit{lectare}）、\textbf{diligent}（勤勉的；\textit{di-}（dis-）分开 + \textit{legere} 选$\to$ 仔细挑选）。delicacy 与 diligent 形近但同源程度有限——都源自"挑选/愉悦"概念。
\end{origin}

% ---------------- 选收的未标 * 词条 ----------------

\vw{glacier}{/ˈɡlæsiə(r)/}{n.}

\sense{n.}{冰川}
\col{a retreating/advancing glacier}
\ex{The glaciers in the Alps have been retreating at an alarming rate.}
\col{glacial lake/valley 冰川湖/冰川谷}
\ex{The turquoise water of glacial lakes attracts tourists from all over the world.}

\begin{origin}
18 世纪借自法语 \textit{glacier}（冰河），法语 \textit{glace}（冰）派生：拉丁语 \textit{glacies}（冰），原始印欧语 *\textit{gel-}（冷、冻）。同源词 \textbf{glacial}（冰河的；冰冷的）、\textbf{glaciate}（使冻结；冰川覆盖）、\textbf{glass}（玻璃——"冰状物"$\to$透明的玻璃）、\textbf{glaze}（上釉；装玻璃）、\textbf{gelid}（冰冷的；\textit{gelid} 是拉丁语直接借词）、\textbf{gelatin}（明胶；"冻"之物）。"冰川"与"玻璃"竟同根——原始印欧人对"冰/透明之物"的认识比表层深得多。
\end{origin}

\vw{suburb}{/ˈsʌbɜːb/}{n.}

\sense{n.}{郊区，城外}
\col{in the suburbs}
\ex{They bought a small house in the suburbs and commute to the city every day.}
\col{a suburb of ...}
\ex{Melbourne is a city of about five million people, with countless suburbs sprawling out from the centre.}
\col{suburban 郊区的}
\ex{He grew up in a quiet suburban neighbourhood in the 1970s.}

\begin{origin}
14 世纪借自拉丁语 \textit{suburbium}（城下之地）：\textit{sub-}（在下）+ \textit{urbs}（城市）。同源词 \textbf{suburban}（郊区的；带贬义时常指"乏味的、缺乏都市气质的"）、\textbf{suburbanite}（郊区居民）、\textbf{suburbia}（郊区整体；常带文化批判意味）、\textbf{urban}（城市的）。suburb 字面"城下"——城墙之下的居所。
\end{origin}

\vw{statue}{/ˈstætʃuː/}{n.}

\sense{n.}{雕像，雕塑}
\col{a bronze/marble statue}
\ex{A bronze statue of the founder stands in the centre of the courtyard.}
\col{the Statue of Liberty 自由女神像}
\ex{The Statue of Liberty was a gift from France to the United States.}

\begin{lex}{statue / statuette / statu 词族与"立"的关联}
\textbf{statue}（雕像，站立姿态）、\textbf{statuette}（小雕像）、\textbf{stature}（身高；声望——"站立的高度"）、\textbf{status}（地位；身份——"站立的位置"）、\textbf{statue}（动产、不动产的法律术语 \textit{statute} 来自拉丁语 \textit{statutum} 已制定之物，但属同根）。这一词族都源自拉丁语 \textit{stare}（站立）的不同派生形式。
\end{lex}

\begin{origin}
14 世纪借自拉丁语 \textit{statua}（雕像、塑像），是 \textit{statutus}（被放置的、被确立的）的名词形式，源于 \textit{statuere}（设立、放置）：\textit{status}（站姿、状态）的动词形式。同源词 \textbf{state}（状态；陈述；国家）、\textbf{statement}（陈述）、\textbf{station}（车站；站）、\textbf{statue}（雕像）、\textbf{static}（静态的）、\textbf{statistics}（统计学）。雕像之所以是"statu"族成员——它"站立着"。
\end{origin}

\vw{camel}{/ˈkæml/}{n.}

\sense{n.}{骆驼}
\col{a camel caravan 骆驼商队}
\ex{Camel caravans have crossed the Sahara for over a thousand years.}
\col{the straw that broke the camel's back 压垮骆驼的最后一根稻草}
\ex{The new tax was the straw that broke the camel's back for many small businesses.}

\begin{origin}
古英语 \textit{camel}，拉丁语 \textit{camelus}，希腊语 \textit{kamēlos}（骆驼），最终可能源自闪族语（希伯来语 \textit{gamal} 骆驼）。同源词 \textbf{camelhair}（驼毛）、法语 \textit{chameau}（骆驼；古怪的人）。"the straw that broke the camel's back" 是英语文化中"最后一击"的经典习语，对应汉语"压垮骆驼的最后一根稻草"。
\end{origin}

\vw{mud}{/mʌd/}{n.}

\sense{n.}{泥，淤泥}
\col{covered in mud}
\ex{After the football match, the players were covered in mud from head to toe.}
\col{mud brick 泥砖}
\ex{Traditional houses in the region are built of mud brick and thatch.}
\col{mud bath 泥浴}
\ex{The spa offers mineral mud baths, said to be good for the skin.}

\begin{gram}{throw / sling mud at sb：抹黑的"泥"隐喻}
\textbf{throw mud at sb}（抹黑、中伤）、\textbf{mud-slinging}（互相揭短）、\textbf{drag sb's name through the mud}（使某人名誉扫地）。mud 在英语政治语言中象征"污浊、污蔑"，是英语"负面物质即负面评价"隐喻的典型案例。
\end{gram}

\begin{origin}
14 世纪中古英语 \textit{mud}，可能与古英语 \textit{mot}（泥土中的一块；与 mote 同源）相关。同源词 \textbf{muddy}（泥泞的）、\textbf{mudslinging}（互相揭短）、荷兰语 \textit{modder}（泥）、德语 \textit{Moder}（腐殖土）。mud 在 14 世纪才进入英语，比 earth、soil、clay 晚得多——古英语多用 mire（淤泥）、slime（粘泥）。
\end{origin}

\vw{brick}{/brɪk/}{n.}

\sense{n.}{砖块}
\col{a brick wall/building}
\ex{The old brick walls of the farmhouse have stood for over two hundred years.}
\col{brick by brick 一砖一瓦地}
\ex{The new library is being built brick by brick, with donations from the local community.}

\sense{n.}{积木}
\col{toy bricks}
\ex{She spent a happy afternoon building castles out of toy bricks.}

\begin{gram}{brick / breeze-block / cinder block / stone 的"块材"细分}
\textbf{brick}（粘土烧制的红砖）、\textbf{breeze-block}（煤渣砖；空心砖）、\textbf{cinder block}（煤渣混凝土砖）、\textbf{stone}（天然石）。英式英语多用 brick（红砖传统），美式英语偏爱 wood frame（木结构）。
\end{gram}

\begin{origin}
古英语 \textit{bryce}（破碎之物；瓦片），与 \textit{brecan}（打破）同源——砖是"打碎泥土后烧成的小块"。同源词 \textbf{break}（打破）、\textbf{breach}（突破）、\textbf{broken}（破碎的）。"砖"与"碎"在古英语中本就是同根词。
\end{origin}

\vw{parcel}{/ˈpɑːsl/}{n. \& vt.}

\sense{n.}{包裹，小包}
\col{a parcel of sth}
\ex{She sent a parcel of books to her brother in Australia.}
\col{parcel post 包裹邮递}
\ex{You can send heavy items by parcel post at a reasonable rate.}

\sense{vt.}{分配（拓展）}
\col{parcel out sth}
\ex{The land was parceled out among the settlers.}
\col{parcel up sth}
\ex{She parcelled up the leftover food and put it in the fridge.}

\begin{lex}{parcel / package / packet / bundle 的"包裹"细分}
\textbf{parcel}（英式"包裹"：邮寄的较大包裹）；\textbf{package}（美式"包裹"；亦作"整套方案"：\textit{a software package} 软件包；\textit{a salary package} 薪资包）；\textbf{packet}（小包、小盒：\textit{a packet of biscuits} 一盒饼干）；\textbf{bundle}（捆扎物：\textit{a bundle of clothes} 一捆衣物）。parcel 与 package 在邮寄语境中常可互换。
\end{lex}

\begin{origin}
14 世纪借自古法语 \textit{parcelle}（一小份），晚期拉丁语 \textit{particella}（小部分），拉丁语 \textit{particula}（极小部分），是 \textit{pars}（部分）的指小形式。同源词 \textbf{part}（部分）、\textbf{particle}（微粒）、\textbf{particular}（特别的）、\textbf{partner}（伙伴）、\textbf{partition}（分隔）、\textbf{parlance}（说法、用语）。"包裹"原义是"分配出的一小份"。
\end{origin}

\vw{parking}{/ˈpɑːkɪŋ/}{n.}

\sense{n.}{停车}
\col{parking lot/space 停车场/停车位}
\ex{The hotel has free parking for guests.}
\col{no parking 禁止停车}
\ex{There is a ``No Parking'' sign at the entrance to the street.}

\sense{n.}{停车位}
\col{find parking 找停车位}
\ex{It took us an hour to find parking in the city centre.}

\begin{lex}{park / parking lot / car park / garage 的停车场所细分}
\textbf{park}（v. 停车；n. 公园）；\textbf{parking lot}（露天停车场，美式通用）；\textbf{car park}（停车场，英式）；\textbf{garage}（车库；可停一至两辆，通常附属于住宅）；\textbf{multistorey car park}（多层停车场）。注意 park 作名词表"公园"与表"停车"同形——不同源。
\end{lex}

\begin{origin}
park（公园）14 世纪借自古法语 \textit{parc}（围起来的场地），最终可能源自日耳曼语 *\textit{parruk}（围栏）。parking（停车）派生自 park（停车），20 世纪由 car park 简缩而来。两种 park 在现代英语中合并为同形词，但词源不同。
\end{origin}

\vw{terrific}{/təˈrɪfɪk/}{adj.}

\sense{adj.}{极好的（口语）}
\ex{We had a terrific time at the beach last weekend.}
\ex{That's a terrific idea — let's do it.}

\sense{adj.}{极大的，可怕的}
\col{a terrific amount/blast}
\ex{The bomb went off with a terrific bang that shattered windows for blocks around.}

\begin{gram}{terrific / terrible / terrified 的褒贬演变}
\textbf{terrific}（极好的；古义"可怕的"）、\textbf{terrible}（糟糕的；古义"可怕的"）、\textbf{terrified}（受到惊吓的）。terrific 与 terrible 同源于拉丁语 \textit{terrere}（恐吓），但现代口语中 terrific 褒义化（类似 great、awesome），terrible 仍表贬义。英语词义褒贬演变中"褒义化"是常见现象。
\end{gram}

\begin{origin}
17 世纪借自拉丁语 \textit{terrificus}（可怕的、引起恐惧的），源于 \textit{terrere}（恐吓）。同源词 \textbf{terrible}（可怕的；糟糕的）、\textbf{terrify}（恐吓）、\textbf{terror}（恐怖）、\textbf{terrorism}（恐怖主义）、\textbf{deter}（威慑；\textit{de-} 离开 + \textit{terrere} $\to$ 吓走）。"极好"与"可怕"在词源中本是一对双生子。
\end{origin}

\vw{delicious}{/dɪˈlɪʃəs/}{adj.}

\sense{adj.}{美味的}
\col{delicious food/wine/cake}
\ex{The chef prepared a delicious meal using only local ingredients.}
\col{smell/taste delicious}
\ex{This soup smells absolutely delicious.}

\sense{adj.}{令人愉快的}
\col{a delicious irony/joke}
\ex{There is a delicious irony in the fact that the two former rivals are now close friends.}

\begin{gram}{delicious / tasty / yummy / scrumptious 的"好吃"语体谱}
\textbf{delicious} 通用，略带正式；\textbf{tasty} 日常口语化；\textbf{yummy} 儿语/非正式；\textbf{scrumptious}（英式口语，"极美味的"）；\textbf{delectable}（书面/文学，\textit{a delectable feast}）。注意 delicious 不可修饰人——说"她很美味"是失礼；可说 \textit{a delicious piece of gossip}（有趣的闲言）。
\end{gram}

\begin{origin}
15 世纪借自古法语 \textit{delicious}，拉丁语 \textit{deliciosus}（使人愉悦的），源于 \textit{delectare}（使人愉悦）。同源词 \textbf{delight}（高兴）、\textbf{delightful}（令人愉快的）、\textbf{delectable}（美味的；令人愉快的）。"美味"与"高兴"在拉丁词源中同根——感官愉悦是认知愉悦的具象。
\end{origin}

\end{document}
